% SpqrBridge slides, the method.
% Build: latexmk -lualatex slides.tex
% Shared preamble of the SpqrBridge slide decks.
% Same style and colors as ref-impl-simulation/slides/methodology.
% Each deck loads it with % Shared preamble of the SpqrBridge slide decks.
% Same style and colors as ref-impl-simulation/slides/methodology.
% Each deck loads it with % Shared preamble of the SpqrBridge slide decks.
% Same style and colors as ref-impl-simulation/slides/methodology.
% Each deck loads it with \input{../common/preamble} and builds with: latexmk -lualatex slides.tex
\documentclass[aspectratio=169,11pt]{beamer}

\usepackage{fontspec}
\setsansfont{DejaVu Sans}
\directlua{
  luaotfload.add_fallback("leanfallback", {"NotoSansMath:mode=harf;", "STIX Two Math:mode=harf;"})
}
\setmonofont{DejaVu Sans Mono}[Scale=0.9, RawFeature={fallback=leanfallback}]

\usepackage{tikz}
\usepackage{array}
\usetikzlibrary{arrows.meta, positioning, calc, fit, backgrounds, shapes.misc, decorations.pathreplacing}

% Colors. Spec side, impl side, proof, trusted, gap.
\definecolor{refc}{HTML}{2B6CB0}
\definecolor{realc}{HTML}{C05621}
\definecolor{proofc}{HTML}{2F855A}
\definecolor{trustc}{HTML}{718096}
\definecolor{gapc}{HTML}{C53030}
\definecolor{ink}{HTML}{1A202C}
\definecolor{paper}{HTML}{FFFFFF}

\setbeamercolor{normal text}{fg=ink, bg=paper}
\setbeamercolor{frametitle}{fg=ink}
\setbeamercolor{title}{fg=ink}
\setbeamercolor{subtitle}{fg=trustc}
\setbeamercolor{structure}{fg=refc}
\setbeamerfont{frametitle}{size=\Large, series=\bfseries}
\setbeamerfont{subtitle}{size=\footnotesize}
\setbeamertemplate{navigation symbols}{}
\setbeamertemplate{footline}{%
  \hfill\textcolor{trustc}{\scriptsize\insertframenumber/\inserttotalframenumber}\hspace{1.2em}\vspace{0.8em}}

\tikzset{
  >={Stealth[length=2.4mm]},
  box/.style={draw, rounded corners=3pt, thick, align=center, inner sep=5pt, font=\small},
  spec/.style={box, draw=refc, fill=refc!8},
  impl/.style={box, draw=realc, fill=realc!8},
  proof/.style={box, draw=proofc, fill=proofc!10},
  trust/.style={box, draw=trustc, fill=trustc!10},
  gap/.style={box, draw=gapc, fill=gapc!6, dashed},
  tool/.style={box, draw=ink!60, fill=ink!4},
  doc/.style={box, draw=ink!60, fill=paper, double, double distance=1pt},
  flow/.style={->, thick, ink!70},
  proved/.style={<->, very thick, proofc},
  rel/.style={-, very thick, proofc, dashed},
  label/.style={font=\scriptsize, align=center},
}

\newcommand{\code}[1]{\texttt{#1}}
\newcommand{\specside}[1]{\textcolor{refc}{#1}}
\newcommand{\implside}[1]{\textcolor{realc}{#1}}

\author{SpqrBridge}
\date{November 2026}
 and builds with: latexmk -lualatex slides.tex
\documentclass[aspectratio=169,11pt]{beamer}

\usepackage{fontspec}
\setsansfont{DejaVu Sans}
\directlua{
  luaotfload.add_fallback("leanfallback", {"NotoSansMath:mode=harf;", "STIX Two Math:mode=harf;"})
}
\setmonofont{DejaVu Sans Mono}[Scale=0.9, RawFeature={fallback=leanfallback}]

\usepackage{tikz}
\usepackage{array}
\usetikzlibrary{arrows.meta, positioning, calc, fit, backgrounds, shapes.misc, decorations.pathreplacing}

% Colors. Spec side, impl side, proof, trusted, gap.
\definecolor{refc}{HTML}{2B6CB0}
\definecolor{realc}{HTML}{C05621}
\definecolor{proofc}{HTML}{2F855A}
\definecolor{trustc}{HTML}{718096}
\definecolor{gapc}{HTML}{C53030}
\definecolor{ink}{HTML}{1A202C}
\definecolor{paper}{HTML}{FFFFFF}

\setbeamercolor{normal text}{fg=ink, bg=paper}
\setbeamercolor{frametitle}{fg=ink}
\setbeamercolor{title}{fg=ink}
\setbeamercolor{subtitle}{fg=trustc}
\setbeamercolor{structure}{fg=refc}
\setbeamerfont{frametitle}{size=\Large, series=\bfseries}
\setbeamerfont{subtitle}{size=\footnotesize}
\setbeamertemplate{navigation symbols}{}
\setbeamertemplate{footline}{%
  \hfill\textcolor{trustc}{\scriptsize\insertframenumber/\inserttotalframenumber}\hspace{1.2em}\vspace{0.8em}}

\tikzset{
  >={Stealth[length=2.4mm]},
  box/.style={draw, rounded corners=3pt, thick, align=center, inner sep=5pt, font=\small},
  spec/.style={box, draw=refc, fill=refc!8},
  impl/.style={box, draw=realc, fill=realc!8},
  proof/.style={box, draw=proofc, fill=proofc!10},
  trust/.style={box, draw=trustc, fill=trustc!10},
  gap/.style={box, draw=gapc, fill=gapc!6, dashed},
  tool/.style={box, draw=ink!60, fill=ink!4},
  doc/.style={box, draw=ink!60, fill=paper, double, double distance=1pt},
  flow/.style={->, thick, ink!70},
  proved/.style={<->, very thick, proofc},
  rel/.style={-, very thick, proofc, dashed},
  label/.style={font=\scriptsize, align=center},
}

\newcommand{\code}[1]{\texttt{#1}}
\newcommand{\specside}[1]{\textcolor{refc}{#1}}
\newcommand{\implside}[1]{\textcolor{realc}{#1}}

\author{SpqrBridge}
\date{November 2026}
 and builds with: latexmk -lualatex slides.tex
\documentclass[aspectratio=169,11pt]{beamer}

\usepackage{fontspec}
\setsansfont{DejaVu Sans}
\directlua{
  luaotfload.add_fallback("leanfallback", {"NotoSansMath:mode=harf;", "STIX Two Math:mode=harf;"})
}
\setmonofont{DejaVu Sans Mono}[Scale=0.9, RawFeature={fallback=leanfallback}]

\usepackage{tikz}
\usepackage{array}
\usetikzlibrary{arrows.meta, positioning, calc, fit, backgrounds, shapes.misc, decorations.pathreplacing}

% Colors. Spec side, impl side, proof, trusted, gap.
\definecolor{refc}{HTML}{2B6CB0}
\definecolor{realc}{HTML}{C05621}
\definecolor{proofc}{HTML}{2F855A}
\definecolor{trustc}{HTML}{718096}
\definecolor{gapc}{HTML}{C53030}
\definecolor{ink}{HTML}{1A202C}
\definecolor{paper}{HTML}{FFFFFF}

\setbeamercolor{normal text}{fg=ink, bg=paper}
\setbeamercolor{frametitle}{fg=ink}
\setbeamercolor{title}{fg=ink}
\setbeamercolor{subtitle}{fg=trustc}
\setbeamercolor{structure}{fg=refc}
\setbeamerfont{frametitle}{size=\Large, series=\bfseries}
\setbeamerfont{subtitle}{size=\footnotesize}
\setbeamertemplate{navigation symbols}{}
\setbeamertemplate{footline}{%
  \hfill\textcolor{trustc}{\scriptsize\insertframenumber/\inserttotalframenumber}\hspace{1.2em}\vspace{0.8em}}

\tikzset{
  >={Stealth[length=2.4mm]},
  box/.style={draw, rounded corners=3pt, thick, align=center, inner sep=5pt, font=\small},
  spec/.style={box, draw=refc, fill=refc!8},
  impl/.style={box, draw=realc, fill=realc!8},
  proof/.style={box, draw=proofc, fill=proofc!10},
  trust/.style={box, draw=trustc, fill=trustc!10},
  gap/.style={box, draw=gapc, fill=gapc!6, dashed},
  tool/.style={box, draw=ink!60, fill=ink!4},
  doc/.style={box, draw=ink!60, fill=paper, double, double distance=1pt},
  flow/.style={->, thick, ink!70},
  proved/.style={<->, very thick, proofc},
  rel/.style={-, very thick, proofc, dashed},
  label/.style={font=\scriptsize, align=center},
}

\newcommand{\code}[1]{\texttt{#1}}
\newcommand{\specside}[1]{\textcolor{refc}{#1}}
\newcommand{\implside}[1]{\textcolor{realc}{#1}}

\author{SpqrBridge}
\date{November 2026}


\title{Transferring secure-messaging's SCKA correctness bound to the extracted SPQR code}
\subtitle{SpqrBridge slides, the method}

\begin{document}

\begin{frame}
  \titlepage
  \centering
  \begin{tikzpicture}[scale=0.8, transform shape, every node/.append style={font=\scriptsize}]
    \node[impl] (c) {\textbf{code}\\extracted \code{States}};
    \node[spec, right=1.3cm of c] (b) {\textbf{\code{Braid}}\\abstract machine};
    \node[spec, right=1.3cm of b] (k) {\textbf{\code{SPQR.Chunked}}\\SM model};
    \node[proof, right=1.3cm of k] (q) {$\Pr[\text{fail}] \le q\varepsilon$\\if four open facts hold};
    \draw[rel] (c) -- node[above, label, text=proofc] {refines} (b);
    \draw[rel] (b) -- node[above, label, text=proofc] {related} (k);
    \draw[flow, proofc] (k) -- node[above, label, text=proofc] {SM proof} (q);
  \end{tikzpicture}
\end{frame}

\begin{frame}[t]{Two proofs about SPQR, one gap}
  \vspace{2.2em}
  \centering
  \begin{tikzpicture}[scale=0.85, transform shape, every node/.append style={font=\scriptsize}]
    \node[impl, text width=5.4cm] (rust) at (0, 2) {Rust code (SparsePostQuantumRatchet)};
    \node[impl, text width=5.4cm] (lean) at (0, 0.6) {Lean code (Aeneas extraction)};
    \node[spec, text width=5.4cm] (braid) at (0, -1.0) {\code{Braid}: ML-KEM Braid,\\with the code's deviations};
    \draw[flow] (rust) -- node[right, label] {Aeneas} (lean);
    \draw[rel] (lean) -- node[right, label, text=proofc, align=left] {successful step $\Rightarrow$ \code{Braid} step\\for \emph{some} randomness} (braid);

    \node[spec, text width=5.4cm] (chunked) at (10.5, 0.6) {\code{SPQR.Chunked}: SM model of the code};
    \node[proof, text width=5.4cm] (bound) at (10.5, -1.0) {correctness bound $1 - \Pr[\text{ok}] \le q\varepsilon$};
    \draw[flow, proofc] (chunked) -- node[right, label, text=proofc] {SM proof} (bound);

    \node[gap, text width=2.4cm] (g) at (5.25, -2.2) {no link between the two};
    \draw[gapc, dashed, thick] (braid.south) |- (g.west);
    \draw[gapc, dashed, thick] (g.east) -| (bound.south);

    \node[label, text=trustc] at (0, 2.9) {SparsePostQuantumRatchet-verify};
    \node[label, text=trustc] at (10.5, 1.5) {secure-messaging (SM)};
  \end{tikzpicture}
\end{frame}

\begin{frame}[t]{A conditional bridge across the gap}
  \vspace{0.6em}
  \centering
  \begin{tikzpicture}[scale=0.8, transform shape, every node/.append style={font=\scriptsize}]
    \node[impl, text width=5.4cm] (rust) at (0, 2.3) {Rust code (SparsePostQuantumRatchet)};
    \node[impl, text width=5.4cm] (lean) at (0, 1.1) {Lean code (Aeneas extraction)};
    \node[spec, text width=5.4cm] (braid) at (0, -1.4) {\code{Braid}: ML-KEM Braid,\\with the code's deviations};
    \draw[flow] (rust) -- node[right, label] {Aeneas} (lean);
    \draw[rel] ([xshift=-1.2cm]lean.south) -- node[left, label, text=proofc, align=right]
      {proved: a successful step\\is a \code{Braid} step\\for \emph{some} randomness} ([xshift=-1.2cm]braid.north);
    \draw[gapc, very thick, dashed] ([xshift=1.2cm]lean.south) -- node[right, label, text=gapc, align=left]
      {assumed (\code{BridgeObligations}), on reachable runs:\\send succeeds, its randomness\\is the RNG's answer, recv is complete} ([xshift=1.2cm]braid.north);

    \node[spec, text width=4.4cm] (chunked) at (10.5, -1.4) {\code{SPQR.Chunked}: SM model\\of the code};
    \node[proof, text width=3.6cm, font=\scriptsize] (bound) at (10.5, -3.1) {SM bound $1 - \Pr[\text{ok}] \le q\varepsilon$};
    \draw[flow, proofc] (chunked) -- node[right, label, text=proofc] {SM proof} (bound);

    \draw[rel] (braid.east) -- node[above, label, text=proofc] {\code{ChunkedRel}, every arm}
      node[below, label, text=proofc] {this repository} (chunked.west);

    \node[impl, text width=5.4cm, font=\scriptsize] (adapter) at (0, -3.1)
      {\textbf{RNG adapter} \code{codeScheme}: \code{fill\_bytes} is an\\oracle with uniform answers; failure $\mapsto$ \code{none}};
    \draw[flow, realc] (lean.west) -- ++(-0.4, 0) |- (adapter.west);

    \node[proof, text width=2.4cm, font=\scriptsize, line width=1.4pt] (codebound) at (4.9, -3.1)
      {\textbf{on the code:}\\$1 - \Pr[\text{ok}] \le q\varepsilon$};
    \draw[flow, realc] (adapter) -- (codebound);
    \draw[flow, proofc, very thick] (bound.west) -- node[above, label, text=proofc]
      {transfer} (codebound.east);

    \node[label, text=trustc] at (0, 3.1) {SparsePostQuantumRatchet-verify};
    \node[label, text=trustc] at (10.5, -0.6) {secure-messaging (SM)};
  \end{tikzpicture}

  \vspace{0.3em}
  \raggedright\footnotesize
  The bound needs the dashed arrow, not the proved one: ``some randomness'' says nothing about its
  distribution. Issues \#785, \#786 and \#787 own the dashed arrow.
\end{frame}

\begin{frame}[t,fragile]{The result}
  \vspace{0.2em}
  \footnotesize
  Assume the four facts about the code in \code{BridgeObligations}, which are not proved yet. The
  send and receive facts cover reachable runs only. Let
  \code{adv} make at most $q$ send queries, counted over both parties, with $q + 2 < 2^{64}$.\\[0.3em]
  Then the probability that the correctness game fails on the code is at most $q\varepsilon$.
  \vspace{0.2em}
{\scriptsize
\begin{verbatim}
theorem code_correctness_error_le (hB : BridgeObligations)
    (sampleInitKey : ProbComp InitKey)
    (adv : SCKACorrectnessAdversary Message) (q : ℕ)
    (hq : SendQueryBound adv q) (hq64 : q + 2 < 2 ^ 64) :
    1 - Pr{let x ← correctnessExp (codeScheme sampleInitKey) adv}
          [x = true] ≤
      (q : ℝ≥0∞) * codeKem.correctnessError ProbCompRuntime.probComp
\end{verbatim}}
  \begin{itemize}\footnotesize
    \item \code{codeScheme}: the code as an SCKA scheme, with uniform bytes for each send. Under
      \code{IdealRng}, the bound holds for the extracted \code{send} (probability deck).
    \item $\varepsilon$: the correctness error of \code{codeKem}, the KEM of the code's ML-KEM
      models. It has no numerical bound yet.
    \item The adversary delivers unmodified recorded messages, in any order and with replay. No
      injected messages, no liveness.
    \item No \code{sorry}. Axioms beyond Lean's three: HKDF, HMAC, one libcrux workaround.
  \end{itemize}
\end{frame}

\begin{frame}[t]{The chain of relations}\label{fr:chain}
  \vspace{0.9em}
  \centering
  \begin{tikzpicture}[every node/.append style={font=\scriptsize}]
    \node[impl, minimum width=2.6cm] (s) at (0, 0) {code state $s$};
    \node[spec, minimum width=2.6cm] (sg) at (5, 0) {\code{Braid} state $\sigma$};
    \node[spec, minimum width=2.6cm] (st) at (10, 0) {\code{Chunked} state $t$};
    \draw[rel] (s) -- node[above, label, text=proofc] {\code{Refines}\\spqr-verify} (sg);
    \draw[rel] (sg) -- node[above, label, text=proofc] {\code{ChunkedRel}\\this repository} (st);
    \draw[decorate, decoration={brace, mirror, amplitude=6pt}, proofc, thick]
      (s.south west) ++(0,-0.2) -- ($(st.south east)+(0,-0.2)$)
      node[midway, below=8pt, label, text=proofc] {state relation: $RA\;s\;t \iff \exists \sigma,\ \code{Refines}\ s\ \sigma \wedge \code{ChunkedRel}\ \sigma\ t$};
  \end{tikzpicture}
  \vspace{1.2em}

  \begin{itemize}\footnotesize
    \item $RA$: a code state and an SM state are related through \code{Braid}. The code state
      refines some \code{Braid} state $\sigma$ (spqr-verify), and $\sigma$ is related to the SM state
      (\code{ChunkedRel}). Both parties use it.
    \item \code{ChunkedRel}: the same of the eleven states, equal epochs and authenticators, byte
      fields as typed values, equal decoders, and each code encoder's \code{u32} counter is SM's counter modulo $2^{32}$. The header that
      \code{Refines} leaves free is taken from the SM state.
    \item $RM$: \code{Braid}'s message map \code{absMsg}, followed by the chunk map
      $c \mapsto (c.\code{index}, c.\code{data})$ into SM.
    \item \code{send\_chunked} and \code{receive\_chunked} cover every arm of both machines.
  \end{itemize}
\end{frame}

\begin{frame}[t]{The Rust decoder meets SM's decoder law}
  \footnotesize
  SM's machines are generic in a \code{StreamDecoder}. SM's own decoder decodes a set of chunks.
  The Rust decoder skips points as they arrive, so its state depends on the order. SM's proofs
  need one law, \code{Lawful}: adding a chunk of $p$ to a decoder that is honest for $p$ keeps it
  honest, and it then decodes nothing or $p$.
  \vspace{0.2em}

  \centering
  \begin{tikzpicture}[scale=0.65, every node/.append style={font=\scriptsize}]
    \draw[->, ink!50] (0,0) -- (6,0) node[right] {$x$};
    \draw[->, ink!50] (0,0) -- (0,2.4) node[above] {$y$};
    \draw[thick, refc, domain=0.2:5.8, samples=60] plot (\x, {1.1 + 0.6*sin(\x*65)});
    \foreach \x in {0.8, 1.9, 3.1, 4.6} {
      \fill[realc] (\x, {1.1 + 0.6*sin(\x*65)}) circle (2.2pt);
    }
    \node[label, text=refc, anchor=west] at (6.2, 1.6) {codeword polynomial of $p$};
    \node[label, text=realc, anchor=west] at (6.2, 0.8) {stored points};
  \end{tikzpicture}
  \vspace{0.3em}

  \raggedright
  \begin{itemize}\footnotesize
    \item The payload $p$ gives 16 slot polynomials over $\mathrm{GF}(2^{16})$, each of degree
      $< K$.
    \item Honest state: every stored point lies on the slot polynomial of $p$. An honest chunk
      adds a point of the same polynomial, or is skipped.
    \item $K$ points with distinct $x$ fix a polynomial of degree $< K$, so a decoded value is $p$.
    \item \code{polyStreamDecoder\_lawful}: one theorem for the four streams, $K \in \{3, 36, 30, 5\}$.
    \item The law is safety only: a decoder that never decodes meets it too.
  \end{itemize}
\end{frame}

\begin{frame}[t]{Moving the bound to the code}
  \footnotesize
  SM's bound comes from \code{correctness\_error\_le\_of\_potential}, over an invariant
  $\mathit{Inv}_2$ and a potential $V_2$ on its game states. The bridge builds $V_1$ from $V_2$.
  Here $s_1 \sim s_2$ relates party states by $RA$ and recorded messages by $RM$, with equal key
  tables, counters and flag.
  \[
    V_1(s_1) =
    \begin{cases}
      1 & \text{if } s_1 \text{ has failed} \\
      \min\bigl(1,\ \inf\{\, V_2(s_2) \mid s_1 \sim s_2,\ \mathit{Inv}_2\ s_2 \,\}\bigr) & \text{otherwise}
    \end{cases}
  \]
  \centering
  \begin{tikzpicture}[every node/.append style={font=\scriptsize}]
    \node[impl, minimum width=1.4cm] (a) at (0, 0.45) {$s_1$};
    \node[impl, minimum width=1.4cm] (a') at (5, 0.45) {$s_1'$};
    \node[spec, minimum width=1.4cm] (b) at (0, -0.45) {$s_2$};
    \node[spec, minimum width=1.4cm] (b') at (5, -0.45) {$s_2'$};
    \draw[flow, realc] (a) -- node[above, label, text=realc] {query to the code} (a');
    \draw[flow, refc] (b) -- node[below, label, text=refc] {same query to SM} (b');
    \draw[rel] (a) -- node[right, label, text=proofc] {$\sim$} (b);
    \node[label, text=realc, anchor=east] at (-1.2, 0.45) {code's game\\(\code{codeScheme})};
    \node[label, text=refc, anchor=east] at (-1.2, -0.45) {SM's game\\(\code{SPQR.Chunked})};
    \draw[rel] (a') -- node[right=0.75cm, label, text=proofc] {$\sim$, or SM has failed} (b');
  \end{tikzpicture}
  \vspace{0.2em}

  \begin{itemize}\footnotesize
    \item One step, for \emph{every} related $s_2$ with $\mathit{Inv}_2\ s_2$:
      $\ \mathbb{E}[V_1(s_1')] \le \mathbb{E}[V_2(s_2')] \le V_2(s_2) + c$, with $c = \varepsilon$
      for a send and $c = 0$ otherwise. The first $\le$ holds because the states are related again,
      or SM failed and $V_2 \ge 1 \ge V_1$ (slide \ref{fr:oblig}).
    \item The bound survives the infimum and the cap. No related $s_2$ gives $V_1 = 1$. No witness
      is chosen, and the games need not show the adversary equal values
      (\code{correctness\_error\_le\_of\_budgeted\_transfer}).
  \end{itemize}
\end{frame}

\begin{frame}[t]{Why 64-bit epochs never overflow}
  \vspace{0.3em}
  \footnotesize
  An adversary with fewer than $2^{64} - 2$ sends cannot push an epoch to the overflow point.
  \vspace{3.2em}

  \centering
  \begin{tikzpicture}[every node/.append style={font=\scriptsize}]
    \node[impl, text width=3.6cm] (c) at (0,0) {\textbf{code}: epochs are \code{u64},\\\code{recv} fails past $2^{64} - 1$};
    \node[spec, text width=3.6cm] (m) at (7,0) {\textbf{SM model}: epochs are $\mathbb{N}$,\\never fails there};
    \node[gap, text width=2.2cm] (g) at (3.5,0) {differ only\\at the limit};
    \draw[gapc, dashed, thick] (c) -- (g) -- (m);
  \end{tikzpicture}
  \vspace{0.8em}

  \raggedright
  \begin{enumerate}\footnotesize
    \item At the limit the two machines differ, so they cannot be related there.
    \item Every epoch is at most the number of successful sends plus one
      (\code{codeGameInv\_epochs}).
    \item So the step obligations are needed only below a cap $B = 2^{64} - 2$ on the sends: fewer
      than $B$ for a send, at most $B$ for a receive, since deliveries go on after the last send.
      With $q + 2 < 2^{64}$, the adversary's $q$ sends stay below $B$.
  \end{enumerate}
\end{frame}

\begin{frame}[t]{Derived transfer obligations}\label{fr:oblig}
  \vspace{0.9em}
  \centering
  \begin{tikzpicture}[every node/.append style={font=\scriptsize},
      key/.style={proof, text width=1.4cm, minimum height=0.9cm},
      what/.style={box, draw=proofc!60, fill=paper, text width=10.6cm, minimum height=0.9cm,
        align=left, execute at begin node={\hyphenpenalty=10000\exhyphenpenalty=10000}}]
    \node[key] (i) at (0, 0) {\textbf{init}};
    \node[what, right=0.2cm of i] {Every pair of initial states that the code can start in is
      related to a pair that SM can start in.};
    \node[key, below=0.25cm of i] (s) {\textbf{send}};
    \node[what, right=0.2cm of s] {The two sends sample the same number of bytes: 64, 32 or none. They
      are coupled on equal samples, and their results are then related.};
    \node[key, below=0.25cm of s] (r) {\textbf{recv}};
    \node[what, right=0.2cm of r] {If SM accepts a recorded message, the code accepts it too,
      with related results. If SM refuses, nothing is required.};
  \end{tikzpicture}
  \vspace{0.8em}

  \begin{itemize}\footnotesize
    \item Key generation samples 64 bytes and encapsulation samples 32. For these two,
      \code{codeParams} builds SM's operations from the code's own primitives, so both sides use
      the same bytes.
    \item The other sends sample nothing. Both results are fixed and related.
    \item Why nothing is required when SM refuses: SM's game has failed already, and the bound
      counts that failure.
  \end{itemize}
\end{frame}

\begin{frame}[t]{What the theorem trusts}\label{fr:trust}
  \vspace{0.2em}
  \centering
  \begin{tikzpicture}[
      layer/.style={box, minimum width=8.2cm, minimum height=0.5cm, inner sep=3pt, font=\scriptsize},
      why/.style={label, text=trustc, anchor=west, align=left}]
    \node[proof, layer] (thm) at (0, 0) {the theorem, checked by the Lean kernel};
    \node[gap, layer, below=0.12cm of thm] (l1) {premises: \code{BridgeObligations}, four fields};
    \node[trust, layer, below=0.12cm of l1] (l2) {logical axioms: HKDF, HMAC, libcrux workaround};
    \node[trust, layer, below=0.12cm of l2] (l3) {model adequacy: FIPS 203 for libcrux, opaque functions};
    \node[trust, layer, below=0.12cm of l3] (l6) {ideal RNG: \code{IdealRng}, a new generator state per send};
    \node[trust, layer, below=0.12cm of l6] (l4) {the extraction: hand rename for the Aeneas bump};
    \node[tool, layer, below=0.12cm of l4] (l5) {tools: Charon, Aeneas and its library, VCVio, Lean};
    \foreach \n/\t in {l1/{\#785, \#786, \#787}, l2/{\code{\#print axioms}}, l3/{review, testing}, l6/{deployment}, l4/{re-extract}, l5/tools} {
      \node[why] at ([xshift=0.25cm]\n.east) {\t};
    }
  \end{tikzpicture}
  \vspace{0.6em}

  \begin{itemize}\footnotesize
    \item \code{\#print axioms} audits only the axioms layer. The models, the opaque functions
      (SHA3, HMAC, libcrux sizes) and the RNG do not appear in it.
    \item A wrong model does not always break a premise: both sides of the bridge can agree on it.
      So the theorem is about the models, and their adequacy is a separate claim.
    \item $\varepsilon$ is built from these models, so a bad primitive makes the bound weak, not
      false.
  \end{itemize}
\end{frame}

\begin{frame}[t]{Status and next steps}
  \centering
  \scriptsize
  \renewcommand{\arraystretch}{1.25}
  \begin{tabular}{l l l}
    \textbf{piece} & \textbf{where} & \textbf{status} \\
    \hline
    code refines \code{Braid} & spqr-verify & proved, one direction \\
    generic stream decoders & SM & proved, merged \\
    \code{Braid} related to \code{SPQR.Chunked} & this repository & proved, every arm \\
    Rust decoder is lawful & this repository & proved \\
    transfer of the bound & this repository & proved, for $q + 2 < 2^{64}$ \\
    extracted \code{send} with a generator & this repository & proved, under \code{IdealRng} \\
    premise \code{mlkem\_sizes} & spqr-verify, \#786 & open \\
    premises \code{send\_total}, \code{recv\_complete} & spqr-verify, \#785 & open \\
    premise \code{send\_refines\_bytes} & spqr-verify, \#787 & open \\
    a number for $\varepsilon$ & \#788, \#799, a model link & open \\
  \end{tabular}\\[0.4em]
  {\tiny spqr-verify: SparsePostQuantumRatchet-verify. SM: secure-messaging.}
  \vspace{0.3em}

  \normalsize
  \begin{tikzpicture}
    \node[proof, text width=12.5cm, align=flush left, font=\scriptsize]
      {\textbf{Next}
       \begin{itemize}
         \item Build the exact \code{BridgeObligations} record from upstream theorems (\#785,
           \#786, \#787). Closing an issue is not enough if its statement differs.
         \item Bound $\varepsilon$: the \#788 workaround, the \code{sampleNTT} cap (\#799), and a
           proof that the two FIPS 203 models agree (README).
       \end{itemize}};
  \end{tikzpicture}
\end{frame}

\begin{frame}[t]{The probability deck}
  \vspace{0.8em}
  \footnotesize
  This deck gives the method and the result. One more deck, \code{slides/probability}, answers
  one question.
  \vspace{0.8em}

  The extraction is deterministic, and SM bounds a probability. Where do the probabilities come
  from, and what must the code prove about its random bytes?
\end{frame}

\end{document}
